\section{Background}

We consider a standard reward-free Markov decision process (MDP), defined by the tuple \mbox{$\gM = (\gS, \gA, P, \gamma, d_0)$}, where $\gS, \gA$ are the state and action space, $P$ is the transition kernel, $\gamma \in (0,1)$ is the discount factor and $d_0$ is the initial state distribution.
Given the MDP $\gM$, a policy \mbox{$\pi: \gS \rightarrow \Delta\gA$}, mapping states to probabilities over actions, its successor measure \citep{dayan1993improving, blier2021} for any initial state-action pair \mbox{$(s_0, a_0) \in \gS \times \gA$} is the discounted cumulative distribution of visiting future states when starting from state $s_0 \in \gS$, take action $a \in \gA$ and following policy $\pi$ thereafter. Formally, it is defined as
\begin{equation}
\label{eq:succesor_measure}
    M^\pi( X | s_0, a_0 ) = \sum_{t\geq0}\gamma^t P(s_{t+1} \in X | s_0, a_0, \pi) \quad \forall X \in \gS.
\end{equation}
Conveniently, the action-value function for policy $\pi$ and any reward function $r: \gS \rightarrow \sR$, can be expressed as
% The derivation of this equivalence is provided in Appendix~\ref{appendix:successor_q_derivation}.
\begin{equation}
\label{eq:q_from_successor}
Q_r^\pi(s,a) := \mathbb{E} \big[ \sum_{t\geq0}\gamma^t r(s_{t+1}) | s,a,\pi \big] = \int_{s'\in\gS} M^\pi(ds' \mid s,a)\, r(s'),
\end{equation}
which decomposes the Q-function between a reward agnostic measure $M^\pi$, modeling the evolution of the policy in the environment, and the reward function.
There exist several algorithms for learning zero-shot policy agents \citep{eysenbachc, borsauniversal, agarwal2024proto}, but in this work we will focus on the Forward Backward (FB) \citep{touati2021learning,touati2023does} algorithm. See related work and \Cref{appendix:extended_rel_work} for extensive analysis.



% --------------------------------------------------
Given a representation space $\gZ \subseteq \sR^d$ (usually the unit hypersphere of radius $\sqrt{d}$, and a family of policies $\pi_{z\in\gZ}$, parameterized by $z$, the FB algorithm aims to learn a \textit{forward} representation $F^\pi : \mathcal{S} \times \mathcal{A} \times \mathcal{Z} \to \mathcal{Z}$ and a \textit{backward} representation $B : \mathcal{S} \to \mathcal{Z}$
such that the successor measure in \Cref{eq:succesor_measure} admits the low-rank factorization
\begin{equation}
\label{eq:fb_factorization}
M^{\pi_z}(ds' \mid s,a) \approx  F(s,a,z)^\top B(s')\rho(ds') \quad \pi_z(s) = \argmax_a F(s,a,z)^\top z,
\end{equation}
where $\rho$ is the state distribution, $z \in \gZ$ is the task embedding of policy $\pi_z$.
The FB networks are trained to minimize a temporal-difference \textit{contrastive} loss over transitions $\{(s,a,s')\}$, independant future states $\{s^+\}$, and reward embeddings $z\sim \nu$, where $\nu$ is a distribution over $\gZ$ specified later:
\begin{align}
\label{eq:fb_loss}
\mathcal{L}_{\text{FB}}(F,B)
&=
\mathbb{E}_{\substack{
z\sim\nu,\,
(s,a,s')\sim\rho,\\
s^+\sim\rho,\,
a'\sim\pi_z(s')
}}
\Big[
\big(
F(s,a,z)^\top B(s^+)
-
\gamma
\overline{F}(s',a',z)^\top
\overline{B}(s^+)
\big)^2
\Big] \notag \\
&\quad
-
2\,\mathbb{E}_{z\sim\nu,\,(s,a,s')\sim\rho}
\Big[
F(s,a,z)^\top B(s')
\Big],
\end{align}
where $\overline{F}$ and $\overline{B}$ denote target networks and $\rho$ is the dataset distribution. In practice, as per \citet{touati2023does} additional regularization terms are used, see \Cref{app:fb_algo_details} for details. 

The $\argmax$ in \Cref{eq:fb_factorization} is approximated by jointly training an actor network with the following loss:
\begin{equation}
\label{eq:actor_loss}
\mathcal{L}_{\text{actor}}(\pi)=-\mathbb{E}_{z\sim\nu,\; s\sim\rho,\; a\sim\pi_z(s)}
\left[
F(s,a,z)^\top z
\right].
\end{equation}

After training, given a reward function $r$ specified at test time, we can find the task embedding $z_r$ parameterizing the optimal learnt policy for reward $r$, by performing a simple linear regression of $r$ onto the span of the learnt features $B$. For FB, the reward embedding can be computed simply via $z_r = \mathbb{E}_{s\sim\rho}\big[B(s)r(s)\big]$ \citep{touati2021learning}, where the expectation is approximated through sampling from the dataset distribution used for training (or a subset).

If \Cref{eq:fb_factorization} holds, then for any reward function $r$ in the linear span of the features, the policy $\pi_{z_r}$ is guaranteed to be optimal with optimal Q-function $Q^*_r(s,a) = F(s,a,z_r)^\top z_r$ \citep{touati2021learning}. In practice, the quality of the learnt representation depends on the chosen factorization $d$ and the coverage of the data $\rho$ used to train them.  In this work, we focus on the second issue.




The optimality criteria of FB \citep{touati2021learning} shows that it is sufficient to sample latents $z$ from a distribution $\nu$ that has support over the entire hypersphere. Previous works \citep{tirinzoni2025zeroshot} propose biasing a fraction $r$ of the latent samples for exploration and training to allocate more model capacity for goal-reaching tasks. Our work focuses on a novel sampling strategy for efficient exploration.\footnote{Remark: Most of existing works leveraging the FB framework \citep{touati2023does, cetin25, pirotta2024fast} are offline methods, where $\rho$ is the dataset distribution collected by decoupled unsupervised exploration methods. Our work is in the \textit{online} setting instead, where we do not have access to external datasets and we collect data transitions jointly with the learning of the FB algorithm.}

